\section{Notation and implementation guide}
\subsection{Objects, indices and operations}
The index $t$ labels an observation time; $k$ and $j$ label the two regimes or eligible experts as specified locally. A subscript selects an element, so $A_{jk}$ is the entry in row $j$, column $k$ of transition matrix $A$. The regime variable $S_t$ is an integer random variable; $X_t,Y_t,C_t,M_t$ are scalar random variables. Lowercase code arrays store their realized numerical samples. The expression $\sum$ denotes addition over the specified index. A superscript naming a method is a label, not exponentiation; the superscript 2 in squared error is multiplication of a number by itself.

The vectors $\theta$, $\mu$ and $q_t$ each contain two scalar entries. $\theta$ (theta) contains chosen structural coefficients, $\mu$ (mu) contains chosen observation means, and $q_t$ contains calculated state probabilities. $\sigma_M$ (sigma) is a chosen scalar standard deviation. $\beta$ (beta) denotes regression slopes; a hat marks an estimated quantity. $\ell$ (ell) denotes a calculated routing logit. $\tau$ (tau) denotes the defined intervention contrast. The set $\mathcal{I}_t$ denotes information available to the filter; $\mathcal{K}_t$ denotes eligible expert indices. These are different objects and must not be interchanged.

\subsection{Synthetic algorithm before code}
The problem is to recover a changing structural coefficient when both confounding and uncertainty about regimes are possible. A naive pooled regression ignores regime changes; an unadjusted regression also mixes the treatment effect with the effect of a shared cause. State-specific regression can address the first issue only when states are supplied or estimated appropriately. Adjustment can address the second issue only under the specified causal assumptions and availability of the confounder.

The reproducible algorithm therefore generates a known two-state process, creates treatment and outcome samples, fits training slopes within the supplied true states, filters noisy state observations, and evaluates coefficient recovery on held-out observations. Hard routing commits to one state; soft routing averages estimates; oracle routing uses the true test state as a diagnostic baseline. The known truth is used for evaluation, never as an input to the test filter. No neural network is trained in this synthetic experiment.

\subsection{Reading the generator and regression code}
\texttt{default\_rng} creates a random-number generator from a recorded seed. Reusing a seed reproduces the random stream within the stated numerical environment. \texttt{np.zeros(n,dtype=int)} allocates an integer array of length $n$; successive assignments update the stored regime. \texttt{A[state[t-1]]} selects the probability row for the previous state. \texttt{rng.choice} draws the next state using that row. Independent calls to \texttt{rng.normal} create disturbance arrays. Indexing \texttt{theta[state]} selects one coefficient for each stored regime and therefore returns an array with one coefficient per observation.

\texttt{NTRAIN:} means all array positions from the training boundary onward. The emission-reversal scenario creates new test observation noise after reversing the means; it does not change treatment or outcome equations. This preserves paired treatment/outcome data while changing the filter's test inputs.

\texttt{np.column\_stack} constructs a two-dimensional regression design: rows are samples, columns are intercept, treatment and optionally confounder. With $n_k$ training samples in state $k$, the design has shape $n_k\times3$ when adjusted and $n_k\times2$ otherwise. \texttt{lstsq} finds coefficients minimizing the sum of squared residuals. The first returned object is the coefficient vector; its second entry, index 1 in Python's zero-based indexing, is the treatment slope. No matrix inverse is formed explicitly.

\subsection{Reading the filter and evaluation code}
\texttt{q @ A} multiplies a length-two probability vector by a two-by-two transition matrix to predict the next state's probabilities. The likelihood calculation subtracts its largest logarithmic value before exponentiation. Since the states have the same observation standard deviation, their common normal-density factor cancels during normalization. \texttt{q /= q.sum()} updates the probability vector in place by dividing by its sum. The output array has one row per observation and two probability columns; the training rows are removed before test evaluation.

\texttt{q @ beta} multiplies a test-length-by-two matrix by a length-two vector, producing one coefficient estimate per test row. \texttt{argmax(axis=1)} chooses the largest probability's column separately in each row; ties select the first index. \texttt{np.full} fills a test vector with a pooled estimate. \texttt{theta[state[test]]} retrieves the simulator's true test coefficient for scoring only. Squared errors penalize large deviations more than absolute errors.

A list of dictionaries collects one result per scenario, replicate and method. Each dictionary maps a field name to a value. A pandas DataFrame turns these records into a table. Grouping by scenario and method produces summary means and standard deviations; the standard error of the mean is the across-replicate standard deviation divided by the square root of 200. The 1,000 temporally dependent test points within a replicate are not treated as 1,000 independent replications.

\subsection{Reading the evidence-router code}
The \texttt{@dataclass(frozen=True)} decorator defines a record with named fields and prevents ordinary reassignment after creation. Type annotations describe intended input types; they are not a comprehensive runtime schema validator. The brace expression collecting expert names creates a set, which contains each name once. Comparing its size with the input length detects duplicate names. A \texttt{ValueError} stops processing rather than silently overwriting an expert.

The router scans each record through an ordered sequence of tests. \texttt{all} requires every numeric field to be finite; \texttt{isfinite} excludes infinities and missing-number values. \texttt{elif} means a later check runs only if earlier conditions did not match, so the output reports the first rejection reason. A list stores eligible identifier--logit pairs; a dictionary maps rejected identifiers to reasons. If no record remains, the function returns abstention immediately. Otherwise it computes stable exponential weights. The final dictionary contains a status, weights and rejection reasons, and never contains a broker order.

For $K$ evidence records, validation and normalization require $O(K)$ time and storage: work and retained entries grow proportionally to the record count, assuming ordinary dictionary and set operations. With the fixed two-state model and fixed regression column count, the diagnostic cost grows approximately linearly with sequence length and the number of replicates. Storing all replicate-level outcomes uses space proportional to the number of scenario--replicate--method combinations.

\section{Complete reference evidence router}
\lstinputlisting[language=Python]{cmtf/evidence_router.py}
\section{Complete synthetic diagnostic}
\lstinputlisting[language=Python]{cmtf/synthetic_study.py}
